\exercisesection{1}

\begin{enumerate}
\def\labelenumi{\arabic{enumi}.}
\tightlist
\item
  设 K 是域，A 是形式幂级数环 K{[}{[}T{]}{]} 中由 T 的系数为 0 的那些级数组成的子环；A 是一个 Noether 局部整环，但不是整闭的。
\end{enumerate}

\begin{enumerate}
\def\labelenumi{(\alph{enumi})}
\item
  证明 A 的每个非零且异于 A 的主理想都含有唯一的形如 \(T^{n} + \lambda T^{n+1}\) 的元素，其中 \(\lambda \in K\) ，\(n \geqslant 2\) ；此外，这个理想含有所有的幂 \(T^{m}\) ，其中 \(m \geqslant n + 2\) 。由此推出，A 的非主理想恰是诸理想 \(AT^{n} + AT^{n+1}\) ，其中 \(n \geqslant 2\) 。
\item
  证明 A 的每个分式理想都是除性的（证明每个非主理想是两个主理想之交）。描述幺半群 D(A) 的结构。A 的素理想只有极大理想 \(m = AT^{2} + AT^{3}\) 与 (0)。
\end{enumerate}

\begin{enumerate}
\def\labelenumi{\arabic{enumi}.}
\setcounter{enumi}{1}
\tightlist
\item
  \begin{enumerate}
  \def\labelenumii{(\alph{enumii})}
  \tightlist
  \item
    在域 K 上两个未定元的因子分解多项式整环 K{[}X, Y{]} 中，给出两个主理想 a, b 的例子，使 \(a + b\) 不是除性的。
  \end{enumerate}
\end{enumerate}

\begin{enumerate}
\def\labelenumi{(\alph{enumi})}
\setcounter{enumi}{1}
\item
  设 K 是域 \(\mathbf{Q}(\mathbf{X})\) （Q 上一个未定元的有理函数域）的二次扩张，它由对 \(\mathbf{Q}(\mathbf{X})\) 添加多项式 \(Y^{2}-2X\in\mathbf{Q}(\mathbf{X})[\mathbf{Y}]\) 的一个根 y 而得到。设 A 是 K 的由 Z、X 与 y 生成的子环；A 是整闭 Noether 整环。证明 \(p=AX+Ay\) 是高 1 的素理想，但 \(p^{2}\) 不是除性的（证明 \(X^{-1}p^{2}\) 是严格包含 p 的素理想）。此外 \(p.(A:p)\neq A\) 。
\item
  设 A 是整闭 Noether 局部整环，其极大理想 m 的高不是 \(\leqslant1\) ；设 p 是 A 的高 1 素理想；对每个整数 i \textgreater{} 0 ，令 \(a_{i} = p + m^{i}\) ；证明 \(\text{div}\left(\bigcap_{i} a_{i}\right)\) 异于 \(\text{sup}(\text{div}(a_{i}))\) 。
\end{enumerate}

\begin{enumerate}
\def\labelenumi{\arabic{enumi}.}
\setcounter{enumi}{2}
\tightlist
\item
  设 A 是 Noether 整环，p 是高 1 的素理想。
\end{enumerate}

  证明 A: \(p \neq A\) （在 A 的分式域中）。（若 \(c \neq 0\) 是 p 的元素，注意 \(p \in \text{Ass}(A/Ac)\) ，并据此推出 Ac: \(p \neq Ac\) ，依据是第 IV 章，§2，习题 30。）对非 Noether 的整环该结果仍成立吗（考虑一个高 1 的非离散赋值环）？

\begin{enumerate}
\def\labelenumi{\arabic{enumi}.}
\setcounter{enumi}{3}
\item
  设 A 是完全整闭整环，m 是 A 的一个极大理想。证明：或者 m 可逆，或者 m 不是除性的且 A: m = A。（注意，若 m 不可逆，则必有 m. \((\mathbf{A}:\mathfrak{m})^{k} = \mathfrak{m}\) 对每个整数 k \textgreater{} 0 成立。）
\item
  设 A 是整环，K 是它的分式域，U 是 \(K^{*}\) 的由 A 的可逆元组成的子群；把群 \(K^{*}/U = G\) 视为按如下关系有序的群，该关系由关系 \(x^{-1}y \in A\) （在 \(K^{*}\) 中）过渡到商而导出（《代数》第 VI 章，§ 1，第 5 节）。
\end{enumerate}

\begin{enumerate}
\def\labelenumi{(\alph{enumi})}
\item
  为使 \(\mathfrak{a} \neq 0\) （ \(\mathbf{K}\) 的一个分式理想）是除性的，当且仅当 \(\mathfrak{a} - \{0\}\) 在 \(\mathbf{G}\) 中的像是 \(\mathbf{G}\) 中的一个优势集（《代数》第 VI 章，§ 1，习题 30）；对每个分式理想 \(\mathfrak{a} \neq 0\) （ \(\mathbf{K}\) 的），\(\mathfrak{a} - \{0\}\) 的像是包含 \(\mathfrak{a} - \{0\}\) 的像的最小优势集。若 \(\mathbf{A}\) 是 Krull 整环，则 \(\mathbf{G}\) 是同构于形如 \(\mathbf{Z}^{(I)}\) 的有序群的一个定向子群的有序群。
\item
  证明：为使同态 \(w\) ——从 \(G\) 到全序群 \(\Gamma\) ——满足：它在 \(A - \{0\}\) 上对复合映射 \(K^* \to K^*/U \xrightarrow{w} \Gamma\) 的限制是一个赋值，当且仅当 \(w\) 是递增同态。
\item
  设 H 是乘积序群 \(Z \times Z\) 中由这样的有序对 \((s_{1}, s_{2})\) 组成的定向子群：它们满足 \(s_{1} + s_{2} \equiv 0 \pmod{2}\)。证明不存在整环 A，使 \(K^{*}/U\) 是同构于 H 的有序群（由 (b) 推出这样的整环必然是 Krull 整环，但对这个环第 5 节的命题 9 不成立（参看习题 22））。
\end{enumerate}

6. 设 A 是一个整环。\par

\begin{enumerate}
\def\labelenumi{(\alph{enumi})}
\item
  为使除性理想 \(\mathfrak{a}\) （ \(\mathbf{A}\) 中的）满足 \(\operatorname{div}(\mathfrak{a})\) 在 \(\mathbf{D}(\mathbf{A})\) 中可逆，当且仅当 \(\mathfrak{a} \colon \mathfrak{a} = \mathbf{A}\) （参看习题 5 及《代数》第 VI 章，§ 1，习题 30）。特别地，为使 \(\mathfrak{a} \colon \mathfrak{a} = \mathbf{A}\) 对每个除性理想 \(\mathfrak{a}\) （ \(\mathbf{A}\) 的）成立，当且仅当 \(\mathbf{A}\) 是完全整闭整环。
\item
  为使 \(\mathfrak{a} \colon \mathfrak{a} = \mathbf{A}\) 对每个有限生成理想 \(\mathfrak{a} \neq 0\) （ \(\mathbf{A}\) 的）成立，当且仅当 \(\mathbf{A}\) 是整闭的。
\end{enumerate}

\begin{enumerate}
\def\labelenumi{\arabic{enumi}.}
\setcounter{enumi}{6}
\item
  设 K 是域，\((v_{\iota})_{\iota \in I}\) 是 K 上满足第 3 节的条件 \((\mathrm{AK}_{\mathrm{I}})\) 与 \((\mathrm{AK}_{\mathrm{II}})\) 且满足下述条件的离散赋值族：对每个整数 r、每个由有理整数组成的族 \((n_{h})_{1 \leqslant h \leqslant r}\) 以及由 I 的互异元素组成的每个族 \((\iota_{h})_{1 \leqslant h \leqslant r}\) ，存在 \(x \in K\) ，使 \(v_{\iota_{h}}(x) = n_{h}\) 对 \(1 \leqslant h \leqslant r\) 成立，且 \(v_{\iota}(x) \geqslant 0\) 对 \(\iota\) （异于诸 \(\iota_{h}\) ）成立。证明诸 \(v_{\iota}\) 的赋值环之交是一个 Krull 整环 A，其分式域是 K，且 \((v_{\iota})_{\iota \in I}\) 是 A 的本性赋值族。（先证 K 是 A 的分式域，从而 A 是 Krull 整环；再证对一切 \(\iota \in I\) ，A 中由使 \(v_{\iota}(x) > 0\) 的 x 组成的素理想必然高为 1，用反证法论证并利用第 4 节命题 6 的推论 2。）
\item
  设 A 是 Krull 整环；证明对每个集合 I，多项式环 \(A[X_{i}]_{i\in I}\) 是 Krull 整环。（注意，若 B 是 Krull 整环，则 \(B[X_{1},\ldots,X_{n}]\) 的本性赋值在 B 上诱导的赋值就是 B 的本性赋值。）
\end{enumerate}

¶ 9. (a) 设 B 是离散赋值环，A = B{[}{[}X{]}{]} 是 B 上一个未定元的形式幂级数环，C = \(A_{x}\) 是分式环 \(S^{-1}A\)，其中 S 是由诸 \(X^{h}\) （h ≥ 0）组成的乘性子集。设 v 是 B 上的一个赋范赋值；对 C 的每个形如 f = ∑n≥h \(b_n X^{n}\) ≠ 0（\(b_h\) 在 B 中 ≠ 0，h 正或负）的元素，写 s(f) = v(\(b_h\) )。证明 s 是 C 上的欧几里得 stathm（《代数》第 VII 章，§ 1，习题 7），且对 C 中非零的 f, g 有 s(fg) = s(f) + s(g)。（若 s(f) = p ≥ s(g) = q 且 h 是 f 中 ≠0 项的诸次数中最小者，证明存在 u ∈ C 使 f - ug = \(X^{h+1} f_1\) ，其中 \(f_1\) ∈ C，并用反证法推出 C 中存在一种 ``欧几里得除法'' 过程。）若 s(f) = 0，则 f 在 C 中可逆。

* 证明 A 不是 Dedekind 整环。 *

\begin{enumerate}
\def\labelenumi{(\alph{enumi})}
\setcounter{enumi}{1}
\tightlist
\item
  由 (a) 推出：若 B 是 Krull 整环，则 A = B{[}{[}X{]}{]} 也是 Krull 整环。（若 K 是 B 的分式域，注意 A 是 K{[}{[}X{]}{]} 与一族 \((\mathrm{C}_{t})\) 与 A 有相同分式域的主理想整环之交，且 A 的每个元素在几乎所有 \(C_{t}\) 中可逆。）
\end{enumerate}

\begin{enumerate}
\def\labelenumi{\arabic{enumi}.}
\setcounter{enumi}{9}
\tightlist
\item
  设 A 是 Krull 整环，K 是它的分式域，L 是包含 K 的域，\((\mathbf{B}_{\alpha})\) 是包含 A 且含于 L 的 Krull 整环的一个右定向族，使得分式域 \(L_{\alpha}\) （ \(B_{\alpha}\) 的）是 K 的有限代数扩张，且 \(B_{\alpha}\) 是 A 在 \(L_{\alpha}\) 中的整闭包。对 A 的每个本性赋值 v 与每个 \(\alpha\) ，令 \(e_{\alpha}(v)\) 为 \(B_{\alpha}\) 上扩张 v 的诸赋值的分歧指数之和。证明：为使诸 \(B_{\alpha}\) 之并是 Krull 整环，当且仅当对 A 的每个本性赋值 v，诸 \(e_{\alpha}(v)\) 组成的集有界。
\end{enumerate}

¶ 11. 除子 \(d\) （在整环 \(\mathbf{A}\) 中）称为有限生成的，如果它形如 \(\operatorname{div}(a)\) ，其中 \(a\) 是一个有限生成分式理想（它未必是除性的；参看 (b)）。

\begin{enumerate}
\def\labelenumi{(\alph{enumi})}
\item
    证明：若 A 是 Krull 整环，则 D(A) 的每个除子都形如 div(a)，其中 \(a = Ax + Ay\) ，x、y 是 A 的分式域中 \(\neq 0\) 的元素（利用第 5 节命题 9 的推论 2）。
\item
  设 K 是域，\(A = K[X_{n}]_{n \in N}\) 是 K 上可数多个未定元的多项式环，它是 Krull 环（习题 8）。设 \(A'\) 是 A 的由单位元与单项式 \(X_{i}X_{j}\) （对所有对 \((i,j)\) ）生成的子环（等价地说，\(A'\) 是所有各项次数均为偶数的多项式组成的集）。若 L 与 \(L'\) 是 A 与 \(A'\) 的分式域，证明 \(A' = A \cap L'\) ，从而 \(A'\) 是 Krull 环。设 p 是 \(A'\) 中由乘积 \(X_{0}X_{i}\) （对所有 \(i \geqslant 0\) ）生成的理想；证明 p 是高 1 的素理想（从而是除性的），但它不是有限生成的。
\end{enumerate}

\begin{enumerate}
\def\labelenumi{\arabic{enumi}.}
\setcounter{enumi}{11}
\tightlist
\item
  \begin{enumerate}
  \def\labelenumii{(\alph{enumii})}
  \tightlist
  \item
    设 \(\mathbf{A}\) 是整闭整环，\(f(\mathbf{X}) = \sum_{i} a_i \mathbf{X}^i\) 与 \(g(\mathbf{X}) = \sum_{j} b_j \mathbf{X}^j\) 是 \(\mathbf{A}[\mathbf{X}]\) 中的两个多项式；证明：若 \(c \in \mathbf{A}\) 使 \(f(\mathbf{X}) g(\mathbf{X})\) 的所有系数都属于 \(\mathbf{A}c\) ，则所有乘积 \(a_i b_j\) 都属于 \(\mathbf{A}c\) （化归为 \(\mathbf{A}\) 是赋值环的情形）。
  \end{enumerate}
\end{enumerate}

\begin{enumerate}
\def\labelenumi{(\alph{enumi})}
\setcounter{enumi}{1}
\item
  若 A 是习题 1 中定义的环，给出两个次数为 1 的多项式 \(f, g\) （在 A{[}X{]} 中）及一个元素 \(c \in \mathbf{A}\) 的例子，使 (a) 的结论对它们不成立（取 \(c = \mathbf{T}^4 + \mathbf{T}^5\) ）。
\item
  在 (a) 的假设下，设 a、b、c 分别是由 f、g 与 fg 的系数生成的 A 的理想：证明 \(\text{div}(c) = \text{div}(a) + \text{div}(b)\) 。给出 \(\tilde{c} \neq \tilde{a}\tilde{b}\) 的一个例子。证明：若 c 是主理想，则 c = ab，且 a 与 b 都可逆；若进而 A 是局部环，则 a 与 b 都是主理想。
\end{enumerate}

\begin{enumerate}
\def\labelenumi{\arabic{enumi}.}
\setcounter{enumi}{12}
\tightlist
\item
  设 A 是整环，\((p_{i})_{i\in I}\) 是一族非零元素，使得诸主理想 \(Ap_{i}\) 都是素理想；设 S 是由诸 \(p_{i}\) 生成的乘性子集。
\end{enumerate}

\begin{enumerate}
\def\labelenumi{(\alph{enumi})}
\item
  证明 \(\mathbf{A} = \mathbf{S}^{-1}\mathbf{A}\cap \left(\bigcap_{i}\mathbf{A}_{\mathbf{A}p_i}\right)\) 。
\item
  假设 A 的主理想的每个非空族都有极大元。证明每个环 \(A_{Ap_{t}}\) 都是离散赋值环（参看第 VI 章，§3，第 5 节，命题 9）。由此推出：若 \(S^{-1}A\) 是整闭的（相应地，完全整闭的），则 A 是整闭的（相应地，完全整闭的）。若 \(S^{-1}A\) 是 Krull 整环，证明 A 是 Krull 整环。
\end{enumerate}

\begin{enumerate}
\def\labelenumi{\arabic{enumi}.}
\setcounter{enumi}{13}
\tightlist
\item
  设 A 是 Krull 整环，S 是 A 的一个不含 0 的饱和乘性子集（《代数》第 II 章，§ 2，习题 1）。证明：若典范同态 \(\bar{i}\colon\mathrm{C}(\mathrm{A})\to\mathrm{C}(\mathrm{S}^{-1}\mathrm{A})\) （第 10 节，命题 17）是双射，则 S 是由
\end{enumerate}
% semantic-fragments: legacy-input-seam
\relax{}

由一族元素 \((p_{\iota})\) 标定，其中每个 \(A p_{\iota}\) 都是素理想（注意，每个与 S 相交的除性素理想都是主理想）。

¶ 15. (a) 设 A 为整环，a、b 为 A 中两个不等于 0 的元素，且 \(Aa \cap Ab = Aab\) 。证明在多项式环 A{[}X{]} 中，由 \(aX + b\) 生成的主理想是素理想（证明每个多项式 \(f(X) \in A[X]\) 在满足 \(f(-b/a) = 0\) 时都是 \(aX + b\) 的倍式，并对 f 的次数作归纳）。

\begin{enumerate}
\def\labelenumi{(\alph{enumi})}
\setcounter{enumi}{1}
\item
  设 A 为 Krull 整环，\(a, b\) 是 A 中的两个元素，且 \(Aa\) 与 \(Aa + Ab\) 是不同的素理想。证明 \(\mathrm{A}[\mathrm{X}] / (a\mathrm{X} + b)\) 是 Krull 整环，且 \(\mathrm{C}(\mathrm{A}[\mathrm{X}] / (a\mathrm{X} + b))\) 同构于 \(\mathrm{C}(\mathrm{A})\)。（借助 (a) 注意到，\(\mathrm{A}[\mathrm{X}] / (a\mathrm{X} + b)\) 同构于 \(\mathrm{A}[b / a] = \mathrm{B}\) ；证明 \(Ba\) 是 B 中的素理想。观察到 \(\mathrm{A}[a^{-1}] = \mathrm{B}[a^{-1}]\) 是 Krull 整环，并使用习题 12。）
\item
  设 A 为整闭 Noether 整环，a、b 为 A 的 Jacobson 根 r 中的两个元素。假设理想 \(Aa + Ab\) 是非除性的素理想；证明 Aa 和 Ab 是素理想。（首先证明 \(Aa \cap Ab = Aab\)，考虑 \(\operatorname{Ass}(A/Aa)\) 中的理想。其次证明关系 \(xy \in Ab\) 和 \(y \notin Aa + Ab\) 对所有 h 都蕴含 \(x \in Ab + Aa^{h}\)，对 h 作归纳；从而得到 \(x \in Ab\)。最后证明，若 \(xy \in Ab\)、\(x \notin Ab\)、\(y \notin Ab\)，则必有 \(x \in Ab + Aa^{h}\) 且 \(y \in Ab + Aa^{h}\) 对所有 h 成立，同样对 h 作归纳；导出矛盾。）
\end{enumerate}

\begin{enumerate}
\def\labelenumi{\arabic{enumi}.}
\setcounter{enumi}{15}
\tightlist
\item
  设 \(A = \bigoplus_{n \in Z} A_n\) 为分次 Krull 整环。令 \(\mathrm{D}_g(\mathrm{A})\)（分别为 \(\mathrm{F}_g(\mathrm{A})\)）表示由分次整除性理想的除子（分别由除子 \(\operatorname{div}(a)\)，其中 a 是 A 中的齐次元）生成的群；记 \(\mathrm{C}_g(\mathrm{A}) = \mathrm{D}_g(\mathrm{A}) / \mathrm{F}_g(\mathrm{A})\) 。令 S 为 A 中的齐次元 \(\neq 0\) 组成的乘性子集。
\end{enumerate}

\begin{enumerate}
\def\labelenumi{(\alph{enumi})}
\item
  证明 A 中每个与 S 相交的高度为 1 的素理想都是分次的（使用第三章 § 1 第 4 节命题 5）。
\item
  令 \(\mathbf{B} = \mathbf{S}^{-1}\mathbf{A}\) ，这是一个分次环（第二章 § 2 第 9 节）；写成 \(\mathbf{B} = \bigoplus_{n\in \mathbf{Z}}\mathbf{B}_n\) ；证明 \(\mathbf{B}_0\) 是域，并且若 \(\mathbf{A}\neq \mathbf{A}_0\) ，则 \(\mathbf{B}\) 同构于环 \(\mathbf{B}_0[\mathbf{X},\mathbf{X}^{-1}]\) ，其元素为 \(\mathrm{P}(\mathbf{X}) / \mathbf{X}^k\) 形式的有理函数，其中 \(\mathrm{P}(\mathbf{X})\in \mathrm{B}_0[\mathbf{X}]\) 。
\item
  证明 \(C_{g}(A)\) 与 \(C(A)\) 典范同构（使用 (a)、(b) 以及第 10 节的公式 (5)）。
\end{enumerate}

\begin{enumerate}
\def\labelenumi{\arabic{enumi}.}
\setcounter{enumi}{16}
\tightlist
\item
  设 \(\mathbf{A} = \bigoplus_{n\in \mathbf{Z}}\mathbf{A}_n\) 为分次 Krull 整环，且 \(p\in \mathbf{A}_1\) 满足 \(\mathbf{A}p\) 是素理想且 \(\neq 0\) 。
\end{enumerate}

\begin{enumerate}
\def\labelenumi{(\alph{enumi})}
\item
  若 B 是第三章 § 1 习题 1 (a) 中定义的环 \(\mathbf{A}_{(p)}\)，证明 B 是 Krull 整环，且群 C(A) 与 C(B) 典范同构。（先证明 C(B) 同构于 C(B{[}p,p \(^{-1}\) {]})，并注意到 B{[}p,p \(^{-1}\) {]} = A{[}p \(^{-1}\) {]}。）
\item
  在习题 15 (b) 的假设下，证明群 \(\mathbf{C}(\mathbf{A})\) 与 \(\mathbf{C}(\mathbf{A}[\mathbf{X},\mathbf{Y}] / (a\mathbf{X} + b\mathbf{Y}))\) 同构。
\end{enumerate}

¶ 18. 设 \(A = \bigoplus_{n \in \mathbf{Z}} A_n\) 为具有正次数的分次 Krull 整环，且 \(A_0\) 是域；令 \(m = \bigoplus_{n \geqslant 1} A_n\) ，它是 \(A\) 的极大理想。令 \(S\) 为齐次元 \(\neq 0\) 的 \(A\) 所成的乘性集，从而环 \(S^{-1}A\) 是主理想整环（习题 16 (b)）。

\begin{enumerate}
\def\labelenumi{(\alph{enumi})}
\item
  设 p 是 A 的一个高度为 1 且与 A --- m 相交的素理想；则 p 不是分次的（否则 p = A），并且 S \(^{-1}\) p 是 S \(^{-1}\) A 的主理想，由形如 x = 1 + x \(_{1}\) + · · · + x \(_{n}\) 的元素生成，其中 x \(_{j}\) ∈ (S \(^{-1}\) A) \(_{j}\)。将 p 中形如 1 + a \(_{1}\) + · · · + a \(_{q}\) 的元素（a \(_{i}\) ∈ A \(_{i}\)）在 S \(^{-1}\) A 中写成 x 的倍数，并使用第五章，§ 1，第 3 节，命题 11，证明 x ∈ p；最后将 p 的每个元素在 S \(^{-1}\) A 中写成 x 的倍数，证明 p = Ax。
\item
 由 (a) 推出 C(A) 与 C(A \(_{m}\) ) 典范同构（使用第 10 节命题 17）。
\end{enumerate}

\begin{enumerate}
\def\labelenumi{\arabic{enumi}.}
\setcounter{enumi}{18}
\tightlist
\item
  设 A 为局部 Krull 整环，m 为其极大理想；若 \(\mathbf{A}' = (\mathbf{A}[\mathbf{X}])_{\mathfrak{m}\mathbf{A}[\mathbf{X}]}\) ，证明 \(\mathbf{C}(\mathbf{A})\) 与 \(\mathbf{C}(\mathbf{A}')\) 典范同构。（使用习题 12 (c)，应用第 10 节命题 17 的判别准则。）
\end{enumerate}

¶ 20. 若整环 A 的每个有限生成理想都是主理想，则称 A 为 Bézout 整环（或 Bezout 整环）。每个 Bézout Noether 整环都是主理想整环。每个赋值环都是 Bézout 整环（因此 Bézout 整环未必是 Noether 环，也未必完全整闭）。

\begin{enumerate}
\def\labelenumi{(\alph{enumi})}
\item
  证明每个 Bézout 整环都是整闭的（参见习题 6）。若一个整环是右定向的 Bézout 子整环族的并，则它是 Bézout 整环。若 A 是 Bézout 整环，则对于 A 的每个乘性子集 S，S \(^{-1}\) A 也是 Bézout 整环，只要 0 \(\notin\) S。
\item
  设 \(v_{i}\)（ \(1 \leqslant i \leqslant n\)）为域 \(K\) 上的独立赋值，\(A_{i}\) 为赋值 \(v_{i}\) 的环。证明 \(A\)（各 \(A_{i}\) 的交）是 Bézout 整环。（若理想 \(a\) 是 \(A\) 的有限生成理想，则 \(v_{i}(x)\)（对 \(x \in a\) 取值）的集合有最小元 \(\alpha_{i}\)，该最小元属于 \(v_{i}\) 的取值群；对所有 \(i\)，取 \(x_{i} \in a\) 使得 \(v_{i}(x_{i}) = \alpha_{i}\) 。使用逼近定理（第六章 §7 第 2 节，
\end{enumerate}

定理 1），证明存在 \(a_i \in \mathbf{A}\)，使得 \(x = \sum_{i=1}^{n} a_i x_i \in \mathfrak{a}\) 满足关系 \(v_i(x) = \alpha_i\)（\(1 \leqslant i \leqslant n\)）。若 \(v_i\) 是离散赋值，则 \(\mathbf{A}\) 是主理想整环。

\begin{enumerate}
\def\labelenumi{(\alph{enumi})}
\setcounter{enumi}{2}
\tightlist
\item
  设 K 为特征 \(\neq2\) 的代数闭域，且 A 为 \(\mathbf{K}(\mathbf{X})\) 的代数闭包中由 \(\mathbf{K}\) 以及两个元素序列 \((x_{n})\)、\((1 / x_{n})\)（ \(1 \leqslant n < +\infty\)）生成的子环，其中 \(x_{1} = \mathbf{X}\) 且 \(x_{n-1} = x_{n}^{2}\)。证明 \(\mathbf{A}\) 是 Bézout 整环（使用 (a)）。若 \(\mathfrak{p}\) 是满足 \(\neq 0\) 的 \(\mathbf{A}\) 中的素理想，证明 \(\mathfrak{p}\) 由形如 \(x_{n} - a_{n}\) 的元素序列生成，其中 \(a_{n} \in \mathbf{K}\)、\(a_{n} \neq 0\) 且 \(a_{n}^{2} = a_{n-1}\)（考虑对所有 \(n\) 的交 \(\mathfrak{p} \cap \mathbf{K}[x_{n}, 1 / x_{n}]\)）。证明 \(\mathfrak{p}\) 不是有限生成的。
\end{enumerate}

\begin{enumerate}
\def\labelenumi{\arabic{enumi}.}
\setcounter{enumi}{20}
\tightlist
\item
  设整环 A 称为伪 Bézout 整环（分别为伪主理想整环），若按习题 5 的记号，群 K*/U 是格序的（完全格序的）。每个 Bézout *（分别为因子分解）* 整环都是伪 Bézout *（分别为伪主理想）* 整环。每个伪主理想整环都是伪 Bézout 整环。一个取值群为 R 的赋值环是伪主理想整环，但不是主理想整环。每个伪 Bézout（分别为伪主理想）整环都是整闭的（分别为完全整闭的）（使用习题 6）。* 每个 Noether 伪 Bézout 整环都是因子分解整环。* 举出一个整闭 Noether（因而为 Krull）整环的例子，它不是伪 Bézout 整环（参见习题 2）。若 A 是伪 Bézout 整环，则对 A 的每个不含 0 的乘性子集 S，S \(^{-1}\) A 也是伪 Bézout 整环（使用第二章 § 2 习题 1）。
\end{enumerate}

¶ 22. (a) 设 \(\Gamma\) 为格序加法群，A 为 \(\Gamma\) 在域 \(k\) 上的代数；A 是整环（代数第二章 §11 第 4 节命题 8）。A 中每个 \(x \neq 0\) 都能唯一写成 \(x = \sum_{i=1}^{n} \alpha_i e^{\nu_i}\)，其中 \(\nu_i\) 是 \(\Gamma\) 中彼此不同的元素，\(e^{\nu_i}\) 是 A 在 \(k\) 上的典范基中相应的元素，而 \(\alpha_i\) 是 \(\neq 0\) 的 \(k\) 中元素。记 \(\phi(x) = \inf(\nu_1, \ldots, \nu_n)\) 于 \(\Gamma\) 中；证明 \(\phi(x + y) \geqslant \inf(\phi(x), \phi(y))\)，若 \(x, y\) 以及 \(x + y\) 在 A 中均为 \(\neq 0\)，且 \(\phi(xy) = \phi(x) + \phi(y)\) 若 \(xy \neq 0\) 在 A 中成立。（为证明第二个断言，先建立下述引理：给定有限族 \((\xi_j)\)（其元素属于 \(\Gamma\)），要使 \(\inf_{i} (\xi_j) = 0\)，充要条件是：对所有 \(\eta > 0\) 的 \(\Gamma\) 中元素，都存在一个指标 \(j\) 和一个元素 \(\zeta \in \Gamma\)，使得 \(0 < \zeta \leqslant \eta\) 且 \(\inf(\xi_j, \zeta) = 0\)。应用此引理，将其归约到 \(\phi(x) = \phi(y) = 0\) 的情形。）

\begin{enumerate}
\def\labelenumi{(\alph{enumi})}
\setcounter{enumi}{1}
\tightlist
\item
  由 (a) 推出，若 \(\mathbf{K}\) 是 \(\mathbf{A},\phi\) 的分式域，则该映射可延拓为从 \(\mathbf{K}^*\) 到 \(\Gamma\) 的同态（仍记为 \(\phi\)），使得
\end{enumerate}

\[
\phi (x + y) \geqslant \inf (\phi (x), \phi (y))
\]

若 \(x, y\) 以及 \(x + y\) 都是 \(\neq 0\) 的 \(\mathbf{K}\) 中元素，则由此推出，若 \(\mathbf{B}\) 是由 \(x \in \mathbf{K}\) 所成且满足 \(x = 0\) 或 \(\phi(x) \geqslant 0\) 的集合，则 \(\mathbf{B}\) 是一个分式域为 \(\mathbf{K}\) 的整环，并且若 \(\mathbf{U}\) 是 \(\mathbf{B}\) 的可逆元群，则 \(\mathbf{K}^*/\mathbf{U}\) 同构于 \(\Gamma\)（特别地，\(\mathbf{B}\) 是伪 Bézout 整环）。导出一个不是完全整闭的伪 Bézout 整环的例子，一个不是伪主理想整环的完全整闭伪 Bézout 整环的例子（参见

代数第六章 § 1 习题 31) * 以及一个不是因子分解整环的伪主理想整环。 *

* (c) 由 (b) 导出一个整环的例子，它是一个右定向的因子分解子整环族的并，完全整闭但不是伪 Bézout 整环。（令 \(\theta\) 为无理数，位于 {[}0, 1{]} 中；对每个整数 \(j\)，令 \(q_j\) 为满足 \(q_j / 2^j < \theta\) 的最大整数。对所有 \(j\)，在乘积群 \(\mathbf{Q} \times \mathbf{Q}\) 上定义格序群结构：取 \((\mathrm{G}_j)_+\) 作为该群中 \(\geqslant 0\) 的元素集，其元素为有序对 \((\xi, \eta)\)，满足 \(\xi \geqslant 0\) 且 \(0 \leqslant \eta \leqslant (q_j 2^{-j}) \xi\)。*

\begin{enumerate}
\def\labelenumi{\arabic{enumi}.}
\setcounter{enumi}{22}
\tightlist
\item
  \begin{enumerate}
  \def\labelenumii{(\alph{enumii})}
  \tightlist
  \item
    设 A 为伪 Bézout 整环（习题 21），K 为其分式域；对每个多项式 \(f \in \mathrm{A}[\mathrm{X}]\)，\(f\) 的系数的一个最大公因子（在 A 的一个可逆元以内确定）称为 \(f\) 的容度。设 \(f, g\) 为 \(\mathrm{A}[\mathrm{X}]\) 中的两个多项式；证明要使 \(f\) 整除 \(g\)（在 \(\mathrm{A}[\mathrm{X}]\) 中），充要条件是 \(f\) 整除 \(g\)（在 \(\mathrm{K}[\mathrm{X}]\) 中），且一个 \(f\) 的容度整除一个 \(g\) 的容度（使用习题 12）（参见习题 30 (c)）。
  \end{enumerate}
\end{enumerate}

\begin{enumerate}
\def\labelenumi{(\alph{enumi})}
\setcounter{enumi}{1}
\item
  由 (a) 推出，若 A 是伪 Bézout（分别为伪主理想）整环，则 A{[}X{]} 也是伪 Bézout（分别为伪主理想）整环。此外，若 \((a_{t})\) 是 A 中不等于 \(\neq0\) 的元素组成的有限族（分别为任意族），且 d 是该族在 A 中的一个最大公因子，则 d 也是该族在 A{[}X{]} 中的一个最大公因子。
\item
  由 (b) 推出，若 A 是伪 Bézout（分别为伪主理想）整环，则 \(A[X_{\lambda}]_{\lambda \in L}\) 是伪 Bézout（分别为伪主理想）整环，其中 \((\mathbf{X}_{\lambda})_{\lambda \in L}\) 是任意族的不定元。
\end{enumerate}

\begin{enumerate}
\def\labelenumi{\arabic{enumi}.}
\setcounter{enumi}{23}
\item
  设 K 为代数闭域，A = K{[}X, Y{]} 为 K 上两个不定元的多项式环，它是 Krull 整环 *（甚至是因子分解整环）。* 对每个有序对 \((\alpha, \beta) \in \mathrm{K}^{2}\)，令 \(w_{\alpha,\beta}\) 为 A 上的离散赋值，使得对于每个 \(f \neq 0\) 的多项式，\(w_{\alpha,\beta}(f)\) 等于 \(\neq 0\) 的单项式在 \(f(\mathrm{X} + \alpha, \mathrm{Y} + \beta)\) 中的最低次数。证明 A 是赋值环 \(w_{\alpha,\beta}\) 的交，且这些赋值中没有一个是 A 的本性赋值。
\item
  若整闭整环 \(\mathbf{A}\) 存在一个赋值族 \((v_{\iota})_{\iota \in I}\)，定义在 \(\mathbf{A}\) 的分式域上并满足性质 \((\mathrm{AK}_{\Pi})\) 和 \((\mathrm{AK}_{\mathrm{III}})\)，则称其具有有限特征。
\end{enumerate}

\begin{enumerate}
\def\labelenumi{(\alph{enumi})}
\tightlist
\item
  证明若 A 整闭且具有有限特征，则对 A 中每个 \(x \neq 0\)，序主除子群中不超过 \(\leqslant \operatorname{div}(x)\) 的极端元只能有有限多个。（令 \(J \subset I\) 为满足 \(v_{\iota}(x) > 0\) 的有限指标集。若 \(m\) 是 J 中元素的数目，证明不可能存在 \(m + 1\) 个元素 \(y_{j}\)（ \(1 \leqslant j \leqslant m + 1\)），它们整除 \(x\) 且 \(\operatorname{div}(y_{j})\) 是极端元；注意对所有 \(j\)，\(\operatorname{div}(y_{j})\) 必须与 \(\sum_{k \neq j} \operatorname{div}(y_{k})\) 互素。）
\end{enumerate}

* (b) 证明第一类超越函数环（第五章 § 1 习题 12）是伪主理想整环（习题 21），但不是有限特征整环（使用 (a)）。 *

¶ 26. 设 A 为整环，K 为其分式域。若 K 上的赋值 v 的赋值环是 A 在 A 的某个素理想 p 处的局部环 \(A_{p}\)（即 A 与 v 的理想的交），则称 v 对 A 是本性赋值。

\begin{enumerate}
\def\labelenumi{(\alph{enumi})}
\item
  设 \(v\) 是对 \(A\) 的高度为 1 的本性赋值，且 \(p\) 是 \(A\) 中由元素 \(x\)（满足 \(v(x) > 0\)）组成的素理想。证明 \(p\) 的高度为 1。若 \(x \in K\) 不属于 \(A_p\)，则 \(A \cap x^{-1}A \subset p\)。若 \(w\) 是 \(K\) 上的赋值，其赋值环 \(B\) 包含 \(A\)，且其理想 \(q\) 满足 \(q \cap A \subset p\)，证明 \(w\) 与 \(v\) 等价。
\item
  若整闭整环 A 的分式域为 K，且存在一个在 K 上、高度为 1 的赋值族 \((v_{\iota})_{\iota\in I}\)，满足性质 \((\mathbf{AK}_{\mathrm{II}})\) 和 \((\mathbf{AK}_{\mathrm{III}})\)，则称 A 具有有限特征且为实型。在这些条件下，证明若 \(z\in K\) 不属于 A，且 p 是满足 \(A\cap z^{-1}A\subset p\) 的 A 的素理想，则存在 \(\iota\in I\)，使 \(v_{\iota}(z)<0\)，并且素理想 \(q_{\iota}\) 由 \(x\in A\) 中满足 \(v_{\iota}(x)>0\) 的元素组成，且包含于 p。（反证法：考虑有限个 \(v_{\iota_{k}}\)，它们满足 \(v_{\iota_{k}}(z)<0\)；证明存在 \(a\in A\)，使 \(a\notin p\) 且对所有 k 都有 \(v_{\iota_{k}}(a)>0\)；推出 \(a^{n}z\notin A\) 对每个整数 n\textgreater0 都成立，并证明这导致矛盾。）由此推出每个对 A 高度为 1 的本性赋值都等价于某个 \(v_{\iota}\)（使用 (a)）。K 的具有有限特征且为实型的子整环的有限交仍是具有有限特征且为实型的整环。
\item
  假设 (b) 的条件成立，并且所有 \(v_{\iota}\) 都是本性赋值。证明对 A 的每个高度为 1 的素理想 p，\(A_{p}\) 都是某个赋值 \(v_{\iota}\) 的环（使用 (b)，取 \(z^{-1} \in p\)）。对所有 \(x \in A\)，主理想 Ax 因而都有唯一的既约准素分解（第四章 § 2 习题 20）；与此分解相应的素理想正是包含 x 的高度为 1 的素理想。
\item
  假设 (b) 的条件成立。令 S 为 A 中不含 0 的乘性子集，且 \(J \subset I\) 是满足 \(\iota \in I\) 且 \(v_{\iota}(x) = 0\)（在 S 上）的指标集合。证明族 \((v_{\iota})_{\iota \in I-J}\) 满足性质 \((\mathrm{AK}_{\mathrm{II}})\) 和 \((\mathrm{AK}_{\mathrm{III}})\)，适用于环 \(S^{-1}A\)；若前述赋值族是 A 的本性赋值族，则 \(S^{-1}A\) 的本性赋值包括 \((v_{\iota})_{\iota \in I}\)。
\item
将第 5 节命题 9 推广到满足 (c) 的假设的情形。
\item
  假设 (b) 的条件成立。令 \(K'\) 为 K 的有限扩张，\(A'\) 为 A 在 \(K'\) 中的整闭包。证明定义在 \(K'\) 上且延拓 \(v_{\iota}\) 的彼此不等价赋值满足 \(AK_{II}\) 和 \(AK_{III}\) 对 \(A'\) 成立，因而后者是具有有限特征且为实型的整环；若 \(v_{\iota}\) 对 A 都是本性赋值，则它们对 \(A'\) 的延拓也是本性赋值（仿照第 8 节命题 12，并使用第六章 § 8 第 3 节注）。
\item
  若 \(\mathbf{A}\) 是具有有限特征且为实型的整环，则 \(\mathbf{A}[\mathbf{X}]\) 也是如此；若 (c) 的条件成立，确定 A{[}X{]} 的本性赋值（仿照第 9 节）。类似地推广习题 8。
\end{enumerate}

\begin{enumerate}
\def\labelenumi{\arabic{enumi}.}
\setcounter{enumi}{26}
\tightlist
\item
  在习题 22 中取 \(\Gamma\) 为族 \((\Gamma_{\iota})_{\iota\in I}\) 的直和，其中 \(\Gamma_{\iota}\) 是加法群 \(\mathbf{R}\) 的子群。证明习题 22 (b) 中定义的环 \(\mathbf{B}\) 是具有有限特征且为实型的整环，并且是 \(\mathbf{B}\) 的一族本性赋值环的交；此外，每个素理想 \(\neq 0\) 的 \(\mathbf{B}\) 都是极大理想。
\end{enumerate}

¶ 28. 若整闭整环 A 的分式域为 K，且存在一个赋值族 \((v_{\iota})_{\iota \in I}\)，它们定义在 K 上，满足性质 \((\mathbf{AK}_{\mathrm{II}})\) 和 \((\mathbf{AK}_{\mathrm{III}})\)，且它们的取值群是加法群 Q 的子群，则称 A 具有有限特征且为有理型。证明 A 的高度为 1 的本性赋值族满足 \((\mathbf{AK}_{\mathrm{II}})\) 和 \((\mathbf{AK}_{\mathrm{III}})\)。（对所有 \(\iota \in I\)，令 \(p_{\iota}\) 为由满足 \(v_{\iota}(x) > 0\) 的 x 组成的 A 的素理想，并令 \(V_{\iota}\) 为赋值 \(v_{\iota}\) 的环。证明若存在两个指标 \(\alpha, \beta\) 使 \(p_{\beta} \subset p_{\alpha}\)，则 A 是 \(V_{\iota}\)（其中指标 \(\iota \neq \alpha\)）的交。为此，反证证明否则存在 \(x \in K^{*}\)、\(y \in p_{\beta}\) 以及两个正整数 r、s，使 \(v_{\beta}(x^{r}y^{s}) > 0\)、\(v_{\alpha}(x^{r}y^{s}) = 0\) 且 \(v_{\iota}(x^{r}y^{s}) \geqslant 0\) 对所有 \(\iota \in I\) 成立，这与假设矛盾。）

\begin{enumerate}
\def\labelenumi{\arabic{enumi}.}
\setcounter{enumi}{28}
\tightlist
\item
  设 A 为整闭整环，K 为其分式域。
\end{enumerate}

\begin{enumerate}
\def\labelenumi{(\alph{enumi})}
\item
  设 L 为 K 的代数扩张，B 为 A 在 L 中的整闭包。证明对 A 的每个分式理想 a，若记 b = aB，则 \(\tilde{b} \cap A = \tilde{a}\)。（将问题化为 \(A \subset \tilde{a}\) 的情形，换言之为 A: \(a \subset A\)，并证明 B: \(b \subset B\)；为此，取 \(x \in B\) : b，令 \(c_i\)（ \(1 \leqslant i \leqslant n\)）为 x 在 K 上的极小多项式的系数；注意对所有 \(y \in a\)，元素 \(c_i y^i\)（ \(1 \leqslant i \leqslant n\)）属于 A，并推出 \(c_i\) 属于 A。）
\item
  令 C 为任意一族不定元上的多项式环 \(\mathbf{A}[\mathbf{X}_{\lambda}]_{\lambda \in L}\)。证明对 A 的每个分式理想 \(\mathfrak{a}\)，若记 \(\mathfrak{c} = \mathfrak{aC}\)，则 \(\tilde{\mathfrak{c}} \cap \mathbf{A} = \tilde{\mathfrak{a}}\)（同法）。
\end{enumerate}

¶ 30. 若在幺半群 D(A) 中，每个有限生成除子（习题 11）都是正则元，则称整环 A 为正则整闭的。每个完全整闭整环都是正则整闭的；每个伪 Bézout 整环（习题 21）都是正则整闭的。

\begin{enumerate}
\def\labelenumi{(\alph{enumi})}
\item
  若 A 是正则整闭的，且 \(d, d', d''\) 是三个有限生成除子，则关系 \(d + d'' \leqslant d' + d''\) 蕴含 \(d \leqslant d'\)。
\item
  由 (a) 推出，整环 A 为正则整闭的充要条件是：对 A 的每个分式理想 a，只要 \(\text{div}(a)\) 是有限生成除子，就有 \(a: a = A\)；特别地，A 是整闭的（习题 6 (b)）。（使用事实 A: (bc) = (A: b): c，其中 b、c 是 A 的两个分式理想，并取 b = a、c = A: a。）此外，\(\text{div}(a)\) 在 D(A) 中是可逆的。
\item
  设 A 为正则整闭整环，K 为其分式域；A 的有限生成除子幺半群 \(\mathrm{D}_{f}(\mathrm{A})\) 在 \(\mathrm{D}(\mathrm{A})\) 中生成一个格序群 \(\mathrm{G}_{f}(\mathrm{A})\)。对每个多项式 \(p \in K[X]\)，令 \(d(p)\) 表示 A 的由 p 的系数生成的分式理想的除子；对 \(\mathrm{K}(\mathrm{X})\) 中的每个有理函数 r = p/q，其中 p、q 属于 \(K[X]\) 且 \(q \neq 0\)，\(d(p) - d(q)\) 在 \(\mathrm{G}_{f}(\mathrm{A})\) 上是良定义的（与将 r 表示为两个多项式之商的方式无关）；记为 \(d(r)\)（参见习题 12 (c)）。若记 \(\gamma(r) = ((r), d(r))\)，其中 (r) 是 \(K[X]\) 中由 r 生成的分式主理想，则 \(\gamma\) 是有序群 \(\mathcal{P}^{*}(\mathrm{A}[\mathrm{X}])\)（§3 第 2 节的记号）到乘积有序群的一个子群的同构
\end{enumerate}

\[
\mathscr {P} ^ {*} (\mathrm{K} [ \mathrm{X} ]) \times \mathrm{G} _ {f} (\mathrm{A}) \subset \mathscr {P} ^ {*} (\mathrm{K} [ \mathrm{X} ]) \times \mathrm{D} (\mathrm{A}).
\]

证明 \(\mathbf{A}[\mathbf{X}]\) 是正则整闭整环，且 \(\mathrm{D}_f(\mathbf{A}[\mathbf{X}])\) 同构于 \(\mathcal{P}^* (\mathbf{K}[\mathbf{X}])\times \mathrm{D}_f(\mathbf{A})\)（使用 (a)、(b) 以及习题 29 (b)）。

\begin{enumerate}
\def\labelenumi{(\alph{enumi})}
\setcounter{enumi}{3}
\tightlist
\item
  设 B 为整闭整环，K 为其分式域，A 为多项式环 K{[}X, Y{]} 的子环，由常数项属于 B 的多项式组成。证明若 \(B \neq K\)，则 A 整闭但不是正则整闭（证明除子 \(\inf(\text{div}(X), \text{div}(Y))\) 对应于某个除性理想 a，使 \(a: a \neq A\)）。
\end{enumerate}

¶ 31. (a) 设 A 为正则整闭整环（习题 30），K 为其分式域；对每个多项式 P ∈ K{[}X{]}，令 c(P) 表示除子 div(a)，其中 a 是 K 中由 P 的系数生成的分式理想。令 B 为 K(X) 的子环，由满足 c(P) ≥ c(Q) 的有理函数 P/Q 组成（使用习题 30 (a) 和 12 (c) 证明此条件只取决于元素 P/Q）。则 B ∩ K = A。

\begin{enumerate}
\def\labelenumi{(\alph{enumi})}
\setcounter{enumi}{1}
\item
  证明 B 是 Bézout 整环（习题 20）（注意，若 P、Q 是 A{[}X{]} 中的两个多项式，则当 m 足够大时，P + X \(^{m}\) Q 在 B 中同时整除 P 和 Q）。
\item
  证明若 A 是 Krull 整环，则 B 是主理想整环（考虑 B 的有限生成分式主理想的递增序列）。举出 A 完全整闭但 B 不是主理想整环的例子（参见习题 22）。
\item
  由 (c) 导出一个 Noether 整环 B 的例子：存在 B 的分式域的一个子域 K，使得 \(B \cap K\) 不是 Noether 环。
\end{enumerate}

\begin{enumerate}
\def\labelenumi{\arabic{enumi}.}
\setcounter{enumi}{31}
\tightlist
\item
  设 A 为整环，\((a_{i})_{1 \leqslant i \leqslant n}\) 为 A 中元素的有限族，a 为理想 \(\sum_{i} A a_{i}\)，R 为 \(A^{n}\) 中由元素 \((a_{1}, \ldots, a_{n})\) 生成的子模。证明 A-模 \(M = A^{n}/R\) 的挠子模成为 M 的直和因子，当且仅当
\end{enumerate}

\[
\mathfrak {a} (\mathrm{A}: \mathfrak {a}) + (\mathrm{A}: (\mathrm{A}: \mathfrak {a})) = \mathrm{A}.
\]

